%! Properties and Applications of Determinants — Unit 5, Lesson 5.2 — Guided Notes
\documentclass[10pt]{article}
\usepackage{linalg-article}
\usepackage{linalg-boxes}
\newcommand{\TallMath}[1]{$\displaystyle #1\rule[-1.4em]{0pt}{3.2em}$}
\newcommand{\vv}{\mathbf{v}}
\newcommand{\ww}{\mathbf{w}}
\newcommand{\avec}{\mathbf{a}}
\newcommand{\bb}{\mathbf{b}}
\newcommand{\xx}{\mathbf{x}}
\newcommand{\T}{\mathsf{T}}
\newcommand{\zero}{\mathbf{0}}

\begin{document}

\pageheader{Unit 5, Lesson 5.2}{Guided Notes}
\namedateperiod

\vspace{0.12in}

\begin{objectivebox}
\textbf{By the end of this lesson, I will be able to\dots}
\begin{itemize}\setlength{\itemsep}{2pt}
  \item state the \emph{rules} determinants obey --- $\det I=1$, a \emph{swap} negates, a \emph{scale} multiplies, and adding a multiple of one row to another changes \emph{nothing};
  \item compute a determinant the \emph{fast} way by reducing to triangular: $\det=\pm(\text{product of the pivots})$;
  \item use the \emph{product rule} $\det(AB)=\det A\,\det B$ and its consequences ($\det A^{-1}=1/\det A$, $\det A=0\Leftrightarrow$ singular).
\end{itemize}
\end{objectivebox}

\vspace{0.06in}

\begin{vocabbox}
\small
Fill in each term as we build it together.
\vspace{2pt}
\termblanklong{Founding rules ($\det I=1$; swap; scale)}
\termblanklong{Shear invariance (add a multiple of a row)}
\termblanklong{Product of pivots (fast determinant)}
\termblanklong{Product rule ($\det AB=\det A\det B$)}
\end{vocabbox}

\begin{hookbox}
Expanding a $3\times3$ by cofactors is \emph{three} $2\times2$ determinants and a lot of bookkeeping.
But you already own a machine that makes any matrix \textbf{triangular} --- elimination (Unit~2) --- and
a triangular determinant is free (Lesson~5.1). The only missing link:
\begin{itemize}\setlength{\itemsep}{1pt}
  \item In the Warm-Up, did the elimination step change the determinant of $M$? \writeline
  \item So if we reduce $A$ all the way to a triangular $U$, how could we read off $\det A$? \writeline
\end{itemize}
\end{hookbox}

\begin{notesbox}{1. The Rules Determinants Live By}
Every rule below is really a statement about the \textbf{volume} of the box the columns span.
\begin{enumerate}[leftmargin=1.5em, itemsep=3pt]
  \item \textbf{The identity.} $\det I=1$ --- the unit cube has volume $1$. \emph{(Our starting point.)}
  \item \textbf{Swap $\Rightarrow$ sign flips.} Exchange two rows (or two columns) and $\det$ is
        \textbf{negated}. Geometrically the box is \emph{reflected}: same size, flipped orientation.
  \item \textbf{Scale $\Rightarrow$ multiply.} Multiply one row by $t$ and $\det$ is multiplied by $t$
        --- stretch one edge by $t$, stretch the volume by $t$.
\end{enumerate}
A quick consequence: if two rows are \emph{equal}, the box is flat, so $\det=$ \blank{1.4cm}.
\end{notesbox}

\begin{notesbox}{2. The Key Property: A Shear Doesn't Change the Volume}
Here is the rule that does all the work: \textbf{adding a multiple of one row to another leaves
$\det$ unchanged.} That move is exactly an \emph{elimination step} (Unit~2), and geometrically it is a
\textbf{shear} --- slide the top of the box sideways. The base and the height never change, so the
\textbf{area (volume) is the same}:

\begin{center}
\begin{tikzpicture}[scale=0.95]
  % original box (upright parallelogram)
  \fill[burgundy!12] (0,0)--(2,0)--(2,1.6)--(0,1.6)--cycle;
  \draw[charcoal, thin] (0,0)--(2,0)--(2,1.6)--(0,1.6)--cycle;
  \draw[<->, linegray] (0,-0.28)--(2,-0.28) node[midway,below,charcoal]{\scriptsize base};
  \node[charcoal] at (1,0.8){\scriptsize area $=bh$};
  % arrow
  \draw[->, very thick, burgundy!70] (2.5,0.8)--(3.5,0.8) node[midway,above]{\scriptsize shear};
  % sheared box (same base, same height, top slid right)
  \fill[greenacc!16] (4,0)--(6,0)--(6.8,1.6)--(4.8,1.6)--cycle;
  \draw[charcoal, thin] (4,0)--(6,0)--(6.8,1.6)--(4.8,1.6)--cycle;
  \draw[<->, linegray] (4,-0.28)--(6,-0.28) node[midway,below,charcoal]{\scriptsize base (same)};
  \draw[dashed, linegray] (6,0)--(6,1.6);
  \draw[<->, linegray] (7.1,0)--(7.1,1.6) node[midway,right,charcoal]{\scriptsize height (same)};
  \node[charcoal] at (5.2,0.8){\scriptsize area $=bh$};
\end{tikzpicture}
\end{center}

\textbf{Check it on a $2\times2$.} Start with $\det\begin{bsmallmatrix} 2 & 1 \\ 1 & 3 \end{bsmallmatrix}=5$.
Now subtract $\tfrac12(\text{column }1)$ from column $2$ (a shear):
\[
  \det\begin{bmatrix} 2 & 1 \\ 1 & 3 \end{bmatrix} = \blank{1.0cm}
  \qquad\text{and}\qquad
  \det\begin{bmatrix} 2 & 0 \\ 1 & 2.5 \end{bmatrix} = \blank{1.0cm}.
\]
Same answer --- the shear slid the box but kept its area.
\end{notesbox}

\begin{notesbox}{3. The Fast Way: Determinant $=$ Product of the Pivots}
Put the rules together. \emph{Shears} (adding a multiple of one row to another) don't change $\det$, and
\emph{swaps} only flip its sign. So run \textbf{elimination} until the matrix is triangular $U$, then
just \textbf{multiply the pivots} --- flipping the sign once for each row swap you made:
\[
  \det A = \pm(\text{product of the pivots}), \qquad (-)\text{ once per row swap}.
\]
\textbf{Worked example (no cofactors).} Reduce
$A=\begin{bsmallmatrix} 1 & 2 & 0 \\ 2 & 5 & 1 \\ 0 & 1 & 3 \end{bsmallmatrix}$:
\[
  \xrightarrow{\;R_2-2R_1\;}
  \begin{bmatrix} 1 & 2 & 0 \\ 0 & 1 & 1 \\ 0 & 1 & 3 \end{bmatrix}
  \xrightarrow{\;R_3-R_2\;}
  \begin{bmatrix} 1 & 2 & 0 \\ 0 & 1 & 1 \\ 0 & 0 & 2 \end{bmatrix}.
\]
No swaps were needed, and the pivots are $1,\;1,\;2$, so
\[
  \det A = (+)\,(1)(1)(2) = \blank{1.2cm}.
\]
That is the same number cofactor expansion would give --- with far less work.
\end{notesbox}

\begin{notesbox}{4. Determinants Multiply: the Product Rule}
When you apply one transformation and then another, the volume factors \emph{multiply}. So determinants
do too:
\[
  \boxed{\;\det(AB) = \det A \,\cdot\, \det B\;}
\]
\textbf{Worked example.} With $\det A=5$ and $\det B=6$,
$\det(AB)=5\cdot 6=\blank{1.0cm}$. Three consequences fall right out:
\begin{itemize}\setlength{\itemsep}{2pt}
  \item \textbf{Inverses:} since $A A^{-1}=I$, we get $\det A\cdot\det A^{-1}=\det I=1$, so
        $\det A^{-1}=\dfrac{1}{\det A}$. (If $\det A=5$, then $\det A^{-1}=$ \blank{1.0cm}.)
  \item \textbf{Singular means zero:} if $\det A=0$ there is \emph{no} number $1/\det A$, so $A$ has no
        inverse --- $\det A=0\Leftrightarrow A$ is \blank{1.9cm}. \emph{(The Unit~2 story again.)}
  \item \textbf{Transpose:} flipping a matrix across its diagonal doesn't change its determinant:
        $\det A^{\T}=\det A$ (so every row rule is also a \emph{column} rule).
\end{itemize}

\vspace{2pt}
\textbf{The road ahead.} \textbf{5.3} \emph{linear transformations}: applying $A$ multiplies every
area/volume by the factor $|\det A|$, and the \emph{sign} of $\det A$ tells you whether orientation was
kept or \blank{1.6cm}.
\end{notesbox}

\begin{practicebox}
Show your work. Use the rules --- reduce to triangular for determinants, and the product rule for the rest.
\begin{enumerate}[leftmargin=1.6em, itemsep=6pt]
  \item Compute $\det\begin{bsmallmatrix} 1 & 3 & 2 \\ 2 & 7 & 5 \\ 0 & 1 & 4 \end{bsmallmatrix}$ by reducing to triangular (product of the pivots). \writeline
  \item You add $3\times(\text{row }1)$ to row $2$ of a matrix. What happens to its determinant? \writeline
  \item If $\det A=4$ and $\det B=2$, find $\det(AB)$ and $\det A^{-1}$. \writeline
\end{enumerate}
\end{practicebox}

\end{document}
